\section{Introduction}

Digital Public Infrastructure is not merely a collection of software stacks; it is the encoding of political arrangements into technical artifacts \citep{winner1980}. Access to essential survival services (food distribution, banking, healthcare, mobility) is increasingly mediated through these systems. Because these systems apply fixed rules to the infinite variety of human life, they will inevitably misclassify, exclude, or fail specific users. This is not a malfunction; it is a mathematical certainty when finite rule sets encounter human diversity.

The critical question for a digital polity is not whether errors occur, but whether the architecture permits those affected to \textbf{detect, correct, and reverse} them before harm becomes permanent.

Current policy discourse defines DPI through aspirational language. The G20 New Delhi Declaration (2023) and the UN/UNDP framework \citep{undp} describe systems that ``should be secure'' or ``can be built on open standards.'' These definitions exclude almost nothing. They articulate intent, not condition. Under such broad definitions, a system that traps users in a coercive loop can be labeled ``public infrastructure'' simply because it operates at scale.

This paper defines corrigibility not as a moral preference, but as a structural property essential for stability. We argue that \textbf{unverifiable constraint is structurally indistinguishable from absolute operator control.} Therefore, public legitimacy requires more than policy assertions; it requires a cryptographic chain of custody and a verifiable capacity for correction.

\paragraph{Paper Structure.} Section~\ref{sec:problem} diagnoses why current DPI definitions fail. Section~\ref{sec:theory} establishes the theoretical foundations from cybernetics, commons governance, and free software. Section~\ref{sec:tests} derives the five structural tests. Section~\ref{sec:dynamics} analyzes temporal dynamics and the correction velocity inequality. Section~\ref{sec:learned} extends the framework to AI systems. Section~\ref{sec:diagnosis} provides empirical evaluation of existing infrastructure. Section~\ref{sec:discussion} addresses implications, including essential services and agentic scaling. The appendices formalize the control-theoretic proofs and provide machine-readable schemas for automated assessment.